\section{连续时间 Markov 链与生成矩阵}
\label{sec:06-Random-Processes/05-Continuous-Time-Markov-Chains-and-Queues/01}

你管理一台关键服务器。运维记录告诉你：平均每 500 小时出一次故障，每次故障平均修 8 小时。你要向上级报告这台机器的可用性——一年 8760 小时里，预期有多少小时它能正常工作。拿计算器一摁：\(500/(500+8) \approx 98.4\%\)，似乎挺高。但你是拿均值算的，背后没模型。万一故障和修复的时间波动很大呢？万一两次故障之间并不独立呢？你需要一个能描述``任何时刻都可能跳变''的随机过程，而不是离散时间 Markov 链那种按固定节拍跳的模型。

设 \(X(t)\) 在离散状态空间 \(S\) 中随连续时间跳变。连续时间意味着状态可以在任意实数时刻改变，不是只在 \(t=1,2,3,\ldots\) 跳。时间齐次 Markov 性说，未来只取决于现在所在的状态，与之前的历史路径以及当前已停留多久都无关：

\[
P(X(t+s)=j\mid X(s)=i,\text{更早历史})
=P(X(t+s)=j\mid X(s)=i).
\]

对每个 \(t\ge 0\)，定义转移矩阵

\[
P(t)_{ij}=P(X(t)=j\mid X(0)=i).
\]

注意这里的 \(t\) 是连续时间——可以是 0.3 小时，也可以是 \(\pi\) 小时，不要求整数。它满足

\[
P(0)=I,
\qquad
P(s+t)=P(s)P(t).
\]

第二式是 Chapman--Kolmogorov 方程。离散时间 Markov 链里它写成 \(P^{(m+n)}=P^{(m)}P^{(n)}\)，这里把上标的步数换成了括号里的连续时间参数，但逻辑完全一样：从 \(i\) 经过时间 \(s+t\) 到 \(j\)，等于先经过时间 \(s\) 到某个中间状态 \(k\)，再经过时间 \(t\) 到 \(j\)，对所有可能的 \(k\) 求和。用 Markov 性把``从 \(k\) 再走 \(t\)''等价于``从 \(k\) 出发走 \(t\)''，就得到矩阵乘法的形式。

\subsection{生成矩阵描述瞬时速率}

离散时间 Markov 链靠一步转移矩阵 \(P\) 描述跳变规律。连续时间没有``一步''这个概念——没有自然的时钟滴答。替代方案是描述``每个瞬间的跳变倾向''，即速率。

取一个极短的时间区间 \(h\downarrow 0\)。对 \(i\ne j\)，定义

\[
q_{ij}
=\lim_{h\downarrow 0}
\frac{P(h)_{ij}}h.
\]

它回答：在状态 \(i\) 的这一瞬间，流向状态 \(j\) 的速率是多少。单位是``每单位时间''——量纲和时间成反比，和概率不同。如果 \(q_{ij}=2\)（每小时），意思是正在状态 \(i\) 时，往 \(j\) 跳的瞬时倾向是每小时 2 次。这不是说一小时之内一定会跳 2 次，而是说在任意瞬间，往 \(j\) 跳的概率近似于 \(2\times h\)（\(h\) 很小）。

把所有离开 \(i\) 的速率加起来，得到总离开速率

\[
q_i=\sum_{j\ne i}q_{ij},
\]

并定义对角元

\[
q_{ii}=-q_i.
\]

矩阵 \(Q=(q_{ij})\) 称为生成矩阵或速率矩阵。它的行不是在描述概率——概率的和应该是 1，而 \(Q\) 每行和为零。对角元是负的：留在 \(i\) 的``速率''是负的，因为停留意味着``没有离开''。满足的性质：

\begin{itemize}
\tightlist
\item
  非对角元 \(q_{ij}\ge 0\)：跳向其他状态的速率不能为负；
\item
  每行和为零：总离开速率等于对角元的绝对值；
\item
  对角元非正：\(q_{ii}\le 0\)，是离开速率的负数。
\end{itemize}

用小时间 \(h\) 内的转移概率近似，可以更直接地感受 \(Q\) 的含义：

\[
P(h)_{ij}=q_{ij}h+o(h)
\quad(i\ne j),
\]
\[
P(h)_{ii}=1-q_ih+o(h).
\]

在极短的时间内，留在当前状态的概率约等于 \(1-q_ih\)；跳向 \(j\) 的概率约等于 \(q_{ij}h\)。\(o(h)\) 是比 \(h\) 更快趋于零的项——它保证了在极短区间内跳两次以上的概率可以忽略。生成矩阵不是转移矩阵：它的行和为零而不是一，对角元可以为负，元素有``每单位时间''量纲。把 \(Q\) 当成``瞬时版的 \(P\)''可以，但别忘了行和不是 1。

\subsection{驻留时间为何 Exponential}

离散时间 Markov 链的驻留时间自动是几何分布——每个时间步以固定概率离开，无记忆性自然成立。连续时间的情况需要推导。

处于状态 \(i\) 后，过程要在某个随机的时刻跳走。问题是：这个等待时间服从什么分布？若过程要保持 Markov 性，已停留多久不能改变剩余驻留时间的分布——否则知道``已经停了 3 小时''和``刚进来''会导致对未来跳走时间的预测不同，那就不是 Markov 的了。

连续时间中唯一具有无记忆性的分布是指数分布。记生存函数 \(S(t)=P(H_i>t)\)，无记忆性要求 \(P(H_i > s+t \mid H_i > s) = P(H_i > t)\)，即 \(S(s+t)=S(s)S(t)\)。这个函数方程在连续且非负解中只有一族解：\(S(t)=e^{-\lambda t}\)。因此当 \(q_i>0\) 时，离开 \(i\) 前的等待时间为

\[
H_i\sim\mathrm{Exponential}(q_i),
\]

其平均为

\[
\E[H_i]=\frac1{q_i}.
\]

离开时跳到 \(j\ne i\) 的条件概率为

\[
r_{ij}=\frac{q_{ij}}{q_i}.
\]

这个比值有个直观解释：每个可能的目的地各有一只速率为 \(q_{ij}\) 的指数时钟在独立竞走，谁的时钟先响就跳向谁。于是 \(j\) 胜出的概率正比于其速率——速率越大的目的地，越有可能先响。赢家通吃，其他时钟在那一刻被丢弃，下次再重新竞走。

于是 CTMC 可以由两部分构造：

\begin{enumerate}
\tightlist
\item
  在状态 \(i\) 停留 \(\mathrm{Exponential}(q_i)\) 时间；
\item
  按概率 \(r_{ij}\) 选择下一状态。
\end{enumerate}

矩阵 \(R=(r_{ij})\) 描述嵌入跳链——它只看跳的顺序和去往哪些状态，不关心跳之间隔了多久。若 \(q_i=0\)，状态 \(i\) 吸收：离开速率为零，驻留时间几乎必然为无穷，此时 \(1/q_i\) 和 \(q_{ij}/q_i\) 都不作为普通有限数使用。

\subsection{跳链平稳分布不等于时间占比}

跳链统计每次跳跃访问的状态频率，CTMC 平稳分布统计连续时间中处于各状态的时间比例。这两个量通常不相等。为什么？因为不同状态的驻留时间不同。

想象一个状态很忙——每次进去很快就出来（\(q_i\) 很大），另一个状态很悠闲——每次进去待很久（\(q_i\) 很小）。如果跳链里两个状态的访问频率相同，时间占比上悠闲状态要远远超过忙碌状态。跳链把每次访问平等计数，时间占比给驻留时间加权。

若跳链平稳分布为 \(\widetilde\pi_i\)，在合适正常返条件下，CTMC 时间平稳概率与

\[
\frac{\widetilde\pi_i}{q_i}
\]

成正比。除以离开速率，就是在补偿平均驻留时间——驻留越久，权重越大。这一关系将在\hyperref[sec:06-Random-Processes/05-Continuous-Time-Markov-Chains-and-Queues/03]{``连续时间平稳分布与流量平衡''}一节用全局平衡方程验证。

\subsection{一个两状态故障---修复模型}

回到开篇那台服务器。状态 \(0\) 为正常，\(1\) 为故障。故障速率 \(\lambda = 1/500\) （每小时），修复速率 \(\mu = 1/8\)（每小时）：

\[
Q=
\begin{bmatrix}
-\lambda&\lambda\\
\mu&-\mu
\end{bmatrix}.
\]

正常平均持续 \(1/\lambda = 500\) 小时，故障平均持续 \(1/\mu = 8\) 小时。跳链每次必在两状态间交替——从正常出去一定是故障，从故障出去一定是正常——跳链的平稳分布是各 1/2。但 CTMC 的时间占比取决于驻留时间：正常状态停 500 小时，故障状态停 8 小时，两个状态的驻留时间差了约 60 倍。于是时间平稳分布大约是 \(98.4\%\) 正常、\(1.6\%\) 故障——开篇用均值心算的结果，现在有了模型支撑。

这个模型在本单元还会反复出现：下一篇求它的瞬态解（``如果在 \(t=0\) 时正常，\(t=100\) 小时之后正常和故障的概率各多少''），``平稳分布''一节算它的时间占比（验证上面的直觉计算）。

这个模型依赖指数寿命与修复时间。指数分布的无记忆性意味着：一台已经正常运行了 400 小时的服务器，和一台刚修好的服务器，下一小时出故障的概率完全相同。如果你知道设备会老化——故障率随着使用时间上升——指数模型就错了。此时仅记录``正常/故障''不足以让状态描述 Markov 化；需要把年龄加进状态，或改用半 Markov 模型、更新过程等更灵活的工具。

\subsection{非爆炸性}

无限状态 CTMC 可能出现一种病态行为：过程在有限时间内完成无穷多次跳跃，跑到无穷远处。这叫爆炸。

为什么会爆炸？因为指数驻留时间的均值是 \(1/q_i\)。如果总离开速率 \(q_i\) 随状态序号增长太快，逐段驻留时间的期望和可能收敛。例如出生率为 \(\lambda_n=n^2\) 的纯出生链，第 \(n\) 个状态平均驻留 \(1/n^2\)，而 \(\sum_{n=1}^\infty 1/n^2 = \pi^2/6\) 收敛——期望意义下有限时间内就跑完了无穷多个状态，路径累积了无限次跳跃。有限状态且速率有限时不会爆炸——状态数有限，速率有界，驻留时间期望的和发散。无限状态中若总速率 \(q_i\) 随状态增长太快，需要额外检查爆炸是否发生。

没有非爆炸性，路径在爆炸时刻之后如何定义会成为问题——爆炸的那一瞬间过程到了哪？转移概率也可能丢失质量：在有限时间内概率质量``泄漏''到了无穷远，\(P(t)\) 的每行之和小于 1。

从生成矩阵的局部速率推导 \(P(t)\) 依赖的微分方程——即 Kolmogorov 前向和后向方程——是下一篇的内容。
